% from ega4-11-fr.tex
\oldpage[IV-3]{122}

Montrons que le diagramme
\begin{equation}
\xymatrix{
&T'\ar[r]^{\partial\otimes1}\ar[d]_{\varphi_1}&
R'\otimes_{A'}N'\ar[r]\ar[d]_{v\otimes1}&
L'_0\otimes_{A'}N'\ar[r]\ar[d]_{u_0\otimes1}&
M'\otimes_{A'}N'\ar[r]\ar@{=}[d]&0\\
0\ar[r]&T''\ar[r]_{\partial''}&R''\otimes_{A'}N'\ar[r]&
L''_0\otimes_{A'}N'\ar[r]&M'\otimes_{A'}N'\ar[r]&0
}
\label{IV.11.2.3.3-fr}
\tag{11.2.3.3}
\end{equation}
est \emph{commutatif}. Pour cela, on vérifie aussitôt (M, IV, 1) que
l'homomorphisme $\partial:T\to R\otimes_AN$ provient (dans le cas actuel)
par passage au quotient de l'homomorphisme
$w:\operatorname{Ker}(f_0\otimes1_N)\to R\otimes_AN$, restriction de
l'homomorphisme $g_0\otimes1_N:L_1\otimes_AN\to R\otimes_AN$, où
$g_0:L_1\to R$ est tel que $f_0=j\circ g_0$~; de même pour $\partial''$.
Il suffit donc de voir que le diagramme
\[
\xymatrix{
\operatorname{Ker}(f_0\otimes1_N)\ar[r]^w\ar[d]&R\otimes_AN\ar[r]&
R'\otimes_{A'}N'=(R\otimes_AN)\otimes_AA'\ar[d]_{v\otimes1}\\
\operatorname{Ker}(f''_0\otimes1_{N'})\ar[rr]&&R''\otimes_{A'}N'
}
\]
est commutatif, ce qui résulte de la commutativité du diagramme
\[
\xymatrix{
L_1\ar[r]\ar[d]&R\ar[d]\\
L''_1\ar[r]&R''.
}
\]

\begin{lemma}[11.2.4]
\label{IV.11.2.4-fr}
Soient $A$ un anneau, $C$ une $A$-algèbre, $M$ un $A$-module, $N$ un
$C$-module, $A'$ une $A$-algèbre. Posons $C'=C\otimes_AA'$,
$M'=M\otimes_AA'$, $N'=N\otimes_AA'=N\otimes_CC'$. Supposons que
$M\otimes_AN$ soit un $C$-module plat. Alors l'homomorphisme canonique
\[
\varphi_1:\operatorname{Tor}_1^A(M,N)\otimes_AA'
\to\operatorname{Tor}_1^{A'}(M',N')
\]
(cf. (11.2.1.5)) est surjectif.
\end{lemma}

Gardons les notations de (11.2.3)~; l'exactitude à droite du produit
tensoriel montre que la suite
$L'_1\xrightarrow{f'_0}L'_0\xrightarrow{\varepsilon'}M'\to0$ est exacte~;
comme $L'_0$ et $L'_1$ sont des $A'$-modules libres, \emph{on peut supposer
que l'on a pris} $L'_0=L''_0$, $L'_1=L''_1$, $u_0$ et $u_1$ étant les
\emph{applications identiques} et $f''_0=f'_0$. Comme
$R=\operatorname{Im}(f_0)$ et $R''=\operatorname{Im}(f''_0)=
\operatorname{Im}(f'_0)$, l'homomorphisme $v$ est surjectif, et il en est
donc de même de $v\otimes1$. La démonstration sera achevée si l'on prouve que
la première ligne de (11.2.3.3) est \emph{exacte}, $v\otimes1$ étant
surjective et $u_0\otimes1$ injectif (Bourbaki, \emph{Alg. comm.},
chap.~I$^{\mathrm{er}}$, \S~1, n$^\circ$~4, cor.~2 de la prop.~2).
Utilisons pour cela

